%% ma: L12-B01-0032
%% lop: 12
%% chuong: 1
%% bai: 1
%% dang: 12.1.21
%% loai: TN4
%% muc: NB
%% dap_an: "D"
%% nguon: "HK1-DeOn_1_1.pdf (M10, Đề ôn tập học kì 1 Toán 12, Đề số 2), Phần I Câu 1"
%% kien_thuc: "Bài 1 (tính đơn điệu và cực trị)"
%% trang_thai: cho-duyet
%% ly_do: "Dạng mới đề xuất 12.1.21"
\begin{ex}
Cho hàm số $y=f(x)$ có đạo hàm trên $\mathbb R$ thoả mãn $f'(x)<0$ $\forall x\in(-1;3)$ và $f'(x)>0$ $\forall x\in(3;5)$. Phát biểu nào sau đây là đúng?
\choice{Hàm số $f(x)$ nghịch biến trên các khoảng $(-1;3)$ và $(3;5)$}{Hàm số $f(x)$ nghịch biến trên khoảng $(3;5)$ và nghịch biến trên khoảng $(-1;3)$}{Hàm số $f(x)$ đồng biến trên các khoảng $(-1;3)$ và $(3;5)$}{\True Hàm số $f(x)$ đồng biến trên khoảng $(3;5)$ và nghịch biến trên khoảng $(-1;3)$}
\loigiai{$f'(x)<0$ trên $(-1;3)$ nên $f$ nghịch biến trên $(-1;3)$; $f'(x)>0$ trên $(3;5)$ nên $f$ đồng biến trên $(3;5)$.}
\end{ex}
